\chapter{Recommendations}

\section{Generic road-obstacle detection}

The camera-based detector described in Chapter~3 is restricted to a predefined set of trained object categories (motor vehicles, motorcycles, bicycles, pedestrians, animals). It cannot, and does not claim to, recognise an arbitrary unexpected object, such as fallen debris, that was not one of its training classes; a class like ``stop sign'' would not fire on a twig lying in the road, since the model is a closed-set classifier over its trained categories, not a generic obstacle detector.

Generic detection of this kind is an active, separate research area, referred to in the literature as road-obstacle or road-anomaly detection, distinct from closed-set object detection. HazardNet addresses the scarcity of real hazard training data by augmenting real road images with synthetic 3D debris models under domain randomisation, allowing detection of debris types not explicitly present in training \cite{choe2023hazardnet}. The SegmentMeIfYouCan benchmark provides datasets and an evaluation protocol specifically for anomalous-object and road-obstacle segmentation, that is, identifying objects that do not belong to the expected road scene at all, rather than classifying among a fixed list \cite{chan2021segmentmeifyoucan}. Open-vocabulary detectors such as Grounding DINO, which pair a Transformer-based detector with language-grounded pre-training to detect objects described by an arbitrary text prompt rather than a fixed class list \cite{liu2023groundingdino}, offer an alternative route, but as large Transformer-based models they are unlikely to run in real time on a Raspberry Pi 4 without a hardware accelerator, which would need to be confirmed by measurement. A more Pi-appropriate direction is a lightweight, custom-trained detector: a 2026 study demonstrates a YOLOv8-n model trained specifically on rare roadside obstacles (cones, debris, barrels, fallen trees, rocks), reporting 98.1\% mAP@0.5 at 68\,fps on GPU hardware \cite{tanveer2026yolov8n}, though this would need re-validation for real-time performance on the Pi 4's more limited compute.

None of these approaches is recommended for adoption within the current project. Each would replace the detection model already implemented, and would require its own dataset, training, and on-hardware evaluation before it could be trusted. Generic road-obstacle detection is therefore recorded here as a concrete, evidenced direction for future work, rather than attempted within the scope or timeline of this dissertation.

Were this pursued in a subsequent project phase, a sensible order of work would be: define the operational meaning of ``unexpected obstacle'' for this vehicle and use case; collect and annotate representative road-hazard imagery, or evaluate an existing benchmark such as RoadObstacle21; compare a custom lightweight detector against an anomaly-segmentation approach on that data; test detection accuracy across daylight, night, rain, shadow, and glare conditions, since the present system has only been characterised under bench lighting; measure inference speed and power draw on the actual Raspberry Pi 4 rather than the GPU hardware the cited studies were evaluated on; calibrate any resulting distance estimate against the VL53L5CX time-of-flight sensor described in Section~\ref{sec:tof-sensor} as an independent reference, rather than trusting a vision-only estimate alone; and only then begin live-camera testing of the new model.
